% ECONOMICS_PROBLEM: {"id": "C72-36", "title": "Exact normal-form correlated equilibrium from an extensive-form tree", "jel_code": "C72", "source_id": "OP-0143"}
% !TeX program = xelatex
\documentclass[11pt,a4paper]{article}
\usepackage[margin=23mm]{geometry}
\usepackage{fontspec}
\setmainfont{Latin Modern Roman}
\usepackage{amsmath,amssymb,mathtools}
\usepackage{xcolor,enumitem,booktabs,longtable,array,xurl,hyperref}
\definecolor{muted}{HTML}{53636C}
\hypersetup{colorlinks=true,urlcolor=blue,linkcolor=blue}
\setlength{\parindent}{0pt}
\setlength{\parskip}{5pt}
\newcommand{\R}{\mathbb R}
\newcommand{\E}{\mathbb E}
\newcommand{\N}{\mathbb N}
\newcommand{\ind}{\mathbf 1}
\newcommand{\Var}{\operatorname{Var}}
\newcommand{\Cov}{\operatorname{Cov}}
\newcommand{\supp}{\operatorname{supp}}
\DeclareMathOperator*{\argmin}{arg\,min}
\DeclareMathOperator*{\argmax}{arg\,max}
\newcommand{\entrymeta}[3]{{\sffamily\small\color{muted}\textbf{\texttt{#1}}\quad|\quad #2\par #3\par}}
\newcommand{\Problemref}[1]{\texttt{#1}}
\newcommand{\JELref}[2]{\texttt{#2} (JEL \texttt{#1})}
\begin{document}
\section*{Exact normal-form correlated equilibrium from an extensive-form tree}
\textbf{Problem ID:} \texttt{C72-36}\quad \textbf{JEL:} \texttt{C72}\par
% BEGIN ECONOMICS STATEMENT
\label{problem:OP-0143}
% GAME_THEORY_PROBLEM: {"local_id": "GT-013", "original_number": 13, "kind": "P", "entry_kind": "statement_only", "published": false, "review_required": true, "category": "game_theory"}
\label{game:GT-013}
{\sffamily\small\color{muted}
\textbf{GT-013} | Correlated equilibrium and representation complexity\\
\textbf{P} | Source-linked exact-computation question; output representation is part of the target.
\par}
\subsection*{Definitions and mathematical statement}
Let $G$ be an explicitly represented finite perfect-recall extensive-form game with rational data. Let $S_i$ be its pure contingent plans: one action at every information set of player $i$. Although one plan has polynomial description length, $|S_i|$ can be exponential in the tree size. For $s\in S=\prod_iS_i$, let $U_i(s)$ be expected payoff over chance.
A normal-form correlated equilibrium (NFCE) is a distribution $\pi$ on $S$ such that
\[
 \sum_{s_{-i}}\pi(s_i,s_{-i})
 \bigl[U_i(s_i,s_{-i})-U_i(s_i',s_{-i})\bigr]\geq0
 \quad\text{for every }i,s_i,s_i'.
\]
The private recommendation is the entire contingent plan $s_i$, disclosed before play. A concrete finite-output question is whether there are a polynomial $p$ and an algorithm which, on every tree of bit length $L$, outputs in $p(L)$ time an exact NFCE as a list
\[
 (s^{1},q_1),\ldots,(s^{K},q_K),\quad
 K\leq p(L),\quad q_\ell\in\mathbb Q_{\geq0},\quad\sum_\ell q_\ell=1,
\]
with total output bit length at most $p(L)$.
\subsection*{Short English statement}
Can one compute an exact correlation over complete plans without expanding the game's strategic form?
\subsection*{Variants and scope boundaries}
For another compact representation, its syntax, sampling/evaluation operations, and allowed output size must be specified. The explicit-support target includes a representation-existence requirement. Extensive-form correlated equilibrium, agent-form correlation, and normal-form coarse correlation use different deviation information. A fixed-error approximate NFCE algorithm does not resolve exact computation.
\subsection*{Source relationship and status}
[S1] explicitly retains Any-NFCE as open and raises representation blowup as a possible obstruction. The polynomial-support output convention here makes that obstruction part of a precise algorithmic statement rather than suppressing it.
\subsection*{References}
\begingroup\footnotesize\raggedright
{\footnotesize\raggedright
\textbf{[S1]} Vincent Cheval, Florian Horn, Soumyajit Paul, and Mahsa Shirmohammadi (2026). \emph{On the Complexity of the Optimal Correlated Equilibria in Extensive-Form Games}. arXiv:2507.11509v2. Locator: Section 2, NFCE definition; Section 7, discussion of Any-NFCE and support-size obstruction. \url{https://arxiv.org/pdf/2507.11509v2}\par}
\par\endgroup

\subsection*{Common conventions: Source mathematical conventions}

\textbf{Finite games.} For a finite set $A$, $\Delta(A)$ is its probability
simplex. Players choose independent mixed strategies in Nash equilibrium;
correlation is allowed only when explicitly part of the solution concept.
In a game with payoff functions $u_i$ and strategy profile $x$, an additive
$\varepsilon$ Nash equilibrium satisfies
\[
 u_i(x)\geq u_i(x_i',x_{-i})-\varepsilon
 \quad\text{for every player $i$ and every admissible deviation $x_i'$}.
\]
For numerical approximation, payoffs are normalized as locally specified.
An $\varepsilon$ payoff residual is distinct from distance at most
$\varepsilon$ to the set of exact equilibria. Point convergence, convergence
of empirical play, and convergence in distance to an equilibrium set are
also distinct.

\textbf{Computational representations.} Unless stated otherwise, a complexity
question takes rational data with binary encoding; polynomial time is measured
in total bit length. A fixed-accuracy polynomial algorithm is not necessarily
polynomial in $1/\varepsilon$, and neither is necessarily polynomial in
$\log(1/\varepsilon)$. Succinct games and value, demand, or sampling oracles
must declare their interfaces. Exact real solutions require an explicit
representation; an unspecified list of arbitrary real numbers is not a
finite computational output. A claimed algorithmic barrier is meaningful
only for its specified model and reductions.

\textbf{Dynamic games.} Histories, information, observation, and allowable
randomization are part of the model. A stationary strategy depends only on
the current state; a positional strategy in a deterministic graph game
chooses one action at each controlled vertex. Nash equilibrium at an initial
state and equilibrium after every history are different requirements.
An expectation of a pathwise $\liminf$ payoff, a $\liminf$ of expected averages,
a limiting expected average, and a uniform finite-horizon equilibrium are
different criteria. None is silently substituted for another.

\textbf{Allocation and mechanisms.} A complete allocation assigns every item
exactly once unless otherwise stated. Zero-valued goods, negative values,
capacity constraints, feasibility, lotteries, and copies of goods affect
fairness definitions and are specified locally. Dominant-strategy incentive
compatibility, Bayesian incentive compatibility, universal truthfulness,
and truthfulness in expectation impose different inequalities. Whenever
an objective ratio has zero denominator, its convention is stated or those
instances are excluded.

\textbf{Infinite-dimensional models.} Expectations, integrals, stochastic
equations, and optimization objectives require their local regularity,
integrability, filtrations, and admissible domains. A backward--forward
mean-field system is a boundary-value problem; it is not an initial-value
semiflow without an additional construction. Long-time, large-population,
and vanishing-noise limits require an order of limits and a convergence
topology. If a minimum need not be attained, the objective is its infimum
or supremum and the approximation requirement is stated separately.
% END ECONOMICS STATEMENT
\end{document}
